% !TEX program = lualatex
\documentclass[10pt,a4paper,twocolumn]{ltjsarticle}

% LuaLaTeX用フォント設定パッケージ
\usepackage{luatexja-fontspec}
\usepackage{amsmath,amssymb}
\usepackage{unicode-math}

% ====================
% 欧文フォント設定 (Libertinus)
% ====================
\setmainfont{Libertinus Serif}[
    BoldFont = {Libertinus Serif Bold},
    ItalicFont = {Libertinus Serif Italic},
    BoldItalicFont = {Libertinus Serif Bold Italic}
]
\setsansfont{Libertinus Sans}[
    BoldFont = {Libertinus Sans Bold},
    ItalicFont = {Libertinus Sans Italic}
]
\setmonofont{DejaVu Sans Mono}[Scale=0.9]

% ====================
% 日本語フォント設定 (原ノ味フォント)
% ====================
\setmainjfont{HaranoAjiMincho-Regular}[
    BoldFont = {HaranoAjiGothic-Medium},
    ItalicFont = {HaranoAjiMincho-Regular},
    BoldItalicFont = {HaranoAjiGothic-Bold}
]
\setsansjfont{HaranoAjiGothic-Regular}[
    BoldFont = {HaranoAjiGothic-Bold}
]
\setmonojfont{HaranoAjiGothic-Regular}

% ====================
% 数式フォント設定 (Libertinus Math)
% ====================
\setmathfont{Libertinus Math}

% ファイル生成日時（JST）
\newcommand{\generatedDate}{2026-01-12}
\newcommand{\generatedTime}{09:50}

% ヘッダー・フッター設定
\usepackage{fancyhdr}
\usepackage{lastpage}
\pagestyle{fancy}
\fancyhf{}
\fancyhead[R]{\small \generatedDate\ \generatedTime\ JST (\thepage/\pageref{LastPage})}
\renewcommand{\headrulewidth}{0.4pt}

% 1ページ目はページ番号なし
\fancypagestyle{firstpage}{
    \fancyhf{}
    \renewcommand{\headrulewidth}{0pt}
}

% その他パッケージ
\usepackage{geometry}
\geometry{margin=18mm}
\usepackage{enumitem}
\setlist{noitemsep,topsep=2pt,parsep=1pt,partopsep=0pt}
\usepackage{booktabs}
\usepackage{array}
\usepackage{tabularx}
\newcolumntype{Y}{>{\raggedright\arraybackslash}X}
\usepackage{ascmac}
\usepackage{fvextra}
\usepackage{hyperref}
\hypersetup{
    colorlinks=true,
    linkcolor=blue,
    urlcolor=blue
}

% セクションスタイル調整
\usepackage{titlesec}
\titleformat{\section}{\normalfont\large\bfseries\sffamily}{\thesection}{1em}{}
\titleformat{\subsection}{\normalfont\normalsize\bfseries\sffamily}{\thesubsection}{1em}{}
\titleformat{\subsubsection}{\normalfont\small\bfseries\sffamily}{\thesubsubsection}{1em}{}
\titlespacing*{\section}{0pt}{2ex plus .5ex minus .2ex}{1ex plus .2ex}
\titlespacing*{\subsection}{0pt}{1.5ex plus .5ex minus .2ex}{0.5ex plus .2ex}

% コードブロック用
\DefineVerbatimEnvironment{code}{Verbatim}{fontsize=\scriptsize,frame=single,framesep=2mm,breaklines=true}

\title{\textsf{\textbf{Media Scribe Workflow 設計原則}}}
\author{massy-Claude Dialogue}
\date{}

\begin{document}

\thispagestyle{firstpage}
\maketitle

\begin{itembox}[l]{本ドキュメントの目的}
media-scribe-workflow の設計判断の根拠となる原則を定義する。下流ドキュメント（機能定義、実装）はこの原則に基づいて設計・評価される。
\end{itembox}

\section{ワークフローの目的}

\subsection{ミッション}

\begin{quote}
多様な動画/音声素材を標準化し、再利用・相互運用可能な形式に編集し、適切な形で文字起こし・レポート化する
\end{quote}

\subsection{最終成果物}

\vspace{0.5\baselineskip}
\noindent{\footnotesize
\begin{tabularx}{\linewidth}{@{}lYY@{}}
\toprule
成果物 & 形式 & 用途 \\
\midrule
構造化動画 & MP4（チャプター付き） & 自己再利用、関係者配布 \\
チャプターファイル & .txt & 外部ツール連携 \\
レポート & PDF & スクリプト・発言記録の文書化 \\
\bottomrule
\end{tabularx}
}
\vspace{0.5\baselineskip}

\subsection{ワークフロー全体像}

\begin{code}
素材（動画/音声/URL）
    ↓
┌────────────────────────┐
│ 構造化フェーズ          │
│ ├─ 素材取得・統合       │
│ ├─ チャプター編集       │
│ │   （人間の判断）      │
│ └─ プロジェクト出力     │
└────────────────────────┘
    ↓
┌────────────────────────┐
│ 出力フェーズ            │
│ ├─ 動画エンコード       │
│ ├─ チャプター分割       │
│ └─ 文字起こし準備       │
└────────────────────────┘
    ↓
┌────────────────────────┐
│ レポートフェーズ        │
│ ├─ 文字起こし           │
│ │  （Whisper/YouTube）  │
│ ├─ AI分析・構造化       │
│ └─ PDF出力              │
└────────────────────────┘
\end{code}

\section{設計原則：配管と陶器の分離}

\subsection{Gitの思想に基づく設計}

\begin{itembox}[l]{設計原則}
「配管（plumbing）がしっかり設計されていれば、どのようなスケールの陶器（porcelain）にも対応できる」
\end{itembox}

\vspace{0.5\baselineskip}
\noindent{\footnotesize
\begin{tabularx}{\linewidth}{@{}lYY@{}}
\toprule
層 & 役割 & 例 \\
\midrule
陶器（Porcelain） & ユーザー向けインターフェース & VCE GUI \\
配管（Plumbing） & 単一目的の処理ツール & vce-encode, vce-split \\
\bottomrule
\end{tabularx}
}
\vspace{0.5\baselineskip}

\subsection{配管優先の原則}

\textbf{開発プロセス:}
\begin{enumerate}
\item 配管ツール（CLI）で機能実装・デバッグ
\item GUIから配管を呼び出し or ロジック共有
\item GUIは「可視化」と「インタラクティブ編集」に専念
\end{enumerate}

\textbf{利点:}
\begin{itemize}
\item デバッグ容易性: CLIは入出力が明確
\item テスト可能性: 自動テストが書きやすい
\item 組み合わせ自由: シェルで柔軟に連携
\item スキル対応: 上級者はCLI、初心者はGUI
\end{itemize}

\subsection{スキルに応じた選択肢}

\textbf{上級ユーザー:}
\begin{code}
for f in *.vce.json; do
    vce-encode "$f" --output "./output/"
done
\end{code}

\textbf{一般ユーザー:}
VCE GUIでプロジェクトを開いてエクスポートボタンを押す。

→ 同じ\texttt{.vce.json}を処理し、同じ結果を得る

\section{プロジェクトファイルの役割}

\subsection{位置づけ}

\begin{itembox}[l]{核心概念}
\texttt{.vce.json} = 人間の判断とマシンの処理を繋ぐ契約
\end{itembox}

\begin{code}
GUI（判断支援）
    │
    ↓
.vce.json ← 判断結果を永続化
    │
    ↓
CLI（処理実行）
\end{code}

\subsection{保存すべき情報}

\vspace{0.5\baselineskip}
\noindent{\footnotesize
\begin{tabularx}{\linewidth}{@{}lY@{}}
\toprule
カテゴリ & 内容 \\
\midrule
ソース情報 & ファイルパス、順序、メタデータ \\
チャプター情報 & 位置、名前、除外フラグ \\
出力設定 & エンコーダ、品質、焼込有無 \\
プロジェクトメタ & 作成日時、バージョン \\
\bottomrule
\end{tabularx}
}
\vspace{0.5\baselineskip}

\subsection{保存しない情報}

\begin{itemize}
\item 波形データ（再生成可能、サイズ大）
\item プレビューキャッシュ（再生成可能）
\item 絶対パス（ポータビリティのため相対パス推奨）
\end{itemize}

\section{UIロックの原則}

\subsection{UIロックとは}

\begin{quote}
UIロック = ユーザーの時間を拘束すること = リアルタイムで注意を要求すること
\end{quote}

\subsection{UIロック不可避な機能}

\textbf{GUIの本質的領域:}

\vspace{0.5\baselineskip}
\noindent{\footnotesize
\begin{tabularx}{\linewidth}{@{}lY@{}}
\toprule
機能 & 理由 \\
\midrule
再生プレビュー & 時間軸メディアの内容確認には実時間が必要 \\
チャプター位置決定 & 「ここで切る」という判断に視聴が必要 \\
除外区間の判断 & 「この部分は不要」の判断に視聴が必要 \\
\bottomrule
\end{tabularx}
}
\vspace{0.5\baselineskip}

\subsection{UIロック回避可能な機能}

\textbf{CLI化すべき領域:}

\begin{itemize}
\item エンコード → バックグラウンド処理 / CLI
\item 分割出力 → CLI
\item 文字起こし → 外部サービス / CLI
\item レポート生成 → CLI（カスタムコマンド）
\end{itemize}

\subsection{設計指針}

\begin{itembox}[l]{GUIの本質的価値}
時間軸上の判断を支援するインターフェース

それ以外の処理は配管に委ねる → ユーザーの時間を不必要にロックしない → バッチ処理・自動化が可能になる
\end{itembox}

\section{人間の判断ポイント}

\subsection{判断の分類}

\textbf{【実時間拘束あり】← GUI必須}
\begin{itemize}
\item チャプター位置
\item 除外区間
\item ソース順序（プレビュー確認時）
\end{itemize}

\textbf{【実時間拘束なし】← 設定UI or 設定ファイルで事前定義可能}
\begin{itemize}
\item 出力形式・品質
\item 文字起こし手段
\item テンプレート選択
\item 詳細度・含める情報
\end{itemize}

\textbf{【一度決めれば再利用可能】← デフォルト化・プリセット化}
\begin{itemize}
\item エンコーダ選択
\item 品質設定
\item タイトル焼込の有無
\item レポートテンプレート
\end{itemize}

\subsection{ワークフロー全体の判断ポイント}

\vspace{0.5\baselineskip}
\noindent{\footnotesize
\begin{tabularx}{\linewidth}{@{}lYl@{}}
\toprule
フェーズ & 判断項目 & 拘束 \\
\midrule
素材取得 & ソース選択、字幕取得要否 & なし \\
構造化 & チャプター位置、除外区間、ソース順序 & \textbf{実時間} \\
出力設定 & 形式、品質、焼込、埋込 & なし \\
文字起こし & 手段、言語、話者識別 & なし \\
レポート & テンプレート、詳細度 & なし \\
\bottomrule
\end{tabularx}
}
\vspace{0.5\baselineskip}

\section{判断支援機能}

\subsection{必要な支援機能}

\vspace{0.5\baselineskip}
\noindent{\footnotesize
\begin{tabularx}{\linewidth}{@{}lYl@{}}
\toprule
判断項目 & 支援機能 & 実装 \\
\midrule
チャプター位置 & 動画プレビュー & ✓ \\
 & 音声再生 & ✓ \\
 & 波形表示 & ✓ \\
 & スペクトログラム & ✓ \\
除外区間 & 除外ハッチング表示 & ✓ \\
ソース順序 & ファイル境界表示 & ✓ \\
\bottomrule
\end{tabularx}
}
\vspace{0.5\baselineskip}

\subsection{判断支援の充足条件}

時間軸メディアに対する人間の判断を支援するために必要なもの:

\begin{enumerate}
\item 知覚手段: 動画プレビュー + 音声再生
\item 全体把握: 波形 + スペクトログラム
\item 結果可視化: 除外ハッチング + ファイル境界表示
\item ナビゲーション: シーク + チャプタージャンプ
\end{enumerate}

→ \textbf{現在の実装で基本的な判断支援は充足}

\section{配管ツール一覧}

\subsection{現在の配管ツール}

\vspace{0.5\baselineskip}
\noindent{\footnotesize
\begin{tabularx}{\linewidth}{@{}lYYl@{}}
\toprule
ツール & 入力 & 出力 & 状態 \\
\midrule
vce-encode & .vce.json & MP4 & 実装済 \\
vce-split & .vce.json & 分割MP4群 & 実装済 \\
vce-chapters & .vce.json & .txt & 未実装 \\
\bottomrule
\end{tabularx}
}
\vspace{0.5\baselineskip}

\subsection{将来の配管ツール（候補）}

\vspace{0.5\baselineskip}
\noindent{\footnotesize
\begin{tabularx}{\linewidth}{@{}lYY@{}}
\toprule
ツール & 入力 → 出力 & 用途 \\
\midrule
vce-init & ファイル群 → .vce.json & プロジェクト初期化 \\
vce-validate & .vce.json → 検証結果 & 整合性チェック \\
vce-report & .vce.json + SRT → レポート入力 & レポート準備 \\
\bottomrule
\end{tabularx}
}
\vspace{0.5\baselineskip}

\section{トレーサビリティ}

\subsection{ドキュメント階層}

\begin{code}
DESIGN_PRINCIPLES.md（本ドキュメント）
    ↓ 参照
vce_feature_matrix.md（機能定義）
    ↓ 参照
実装（コード）
\end{code}

\subsection{設計判断の記録}

設計判断を行う際は、本ドキュメントのどの原則に基づくかを明記する。

\textbf{例:}
\begin{quote}
「エンコード処理をCLI化する」\\
→ 根拠: §4.3 UIロック回避可能な機能、§2.2 配管優先の原則
\end{quote}

\subsection{変更履歴}

\vspace{0.5\baselineskip}
\noindent{\footnotesize
\begin{tabularx}{\linewidth}{@{}lY@{}}
\toprule
日付 & 変更内容 \\
\midrule
2026-01-12 & 初版作成 \\
\bottomrule
\end{tabularx}
}
\vspace{0.5\baselineskip}

\section*{参照ドキュメント}

\begin{itemize}
\item \href{./vce_feature_matrix.md}{vce\_feature\_matrix.md} - VCE機能定義（What/How、因果関係）
\item \href{../CLAUDE.md}{CLAUDE.md} - 開発者向けガイド
\end{itemize}

\section*{所感}

本設計原則は、「配管と陶器の分離」というGitの設計思想を、メディア編集ワークフローに適用したものである。

特筆すべきは「UIロック」という概念の導入である。GUIの価値を「時間軸上の判断支援」に限定し、それ以外の処理をCLI（配管）に委ねるという明確な責任分離は、開発効率とユーザー体験の両面で合理的である。

一方で、この設計が実際に機能するためには、\texttt{.vce.json}のスキーマ設計が極めて重要となる。プロジェクトファイルが「人間の判断とマシンの処理を繋ぐ契約」として機能するためには、十分な表現力と将来の拡張性を両立させる必要がある。

また、「上級者はCLI、初心者はGUI」という二分法は単純化しすぎている可能性もある。実際には、同一ユーザーが状況に応じて両方を使い分けるケースが多いだろう。その意味で、GUIからCLIコマンドを生成・表示する機能や、CLIの実行結果をGUIで確認する機能など、両者の橋渡しとなる機能も検討に値する。

\end{document}
